\documentclass[11pt]{amsart}
\usepackage[T1]{fontenc}
\usepackage{lmodern,amsmath,amssymb,mathtools,microtype,booktabs}
\usepackage[margin=1.15in]{geometry}
\usepackage[hidelinks]{hyperref}
\newtheorem{theorem}{Theorem}[section]
\newtheorem{lemma}[theorem]{Lemma}
\newtheorem{proposition}[theorem]{Proposition}
\newtheorem{corollary}[theorem]{Corollary}
\theoremstyle{definition}\newtheorem{definition}[theorem]{Definition}
\theoremstyle{remark}\newtheorem{remark}[theorem]{Remark}
\newcommand{\R}{\mathbb R}\newcommand{\Z}{\mathbb Z}
\newcommand{\E}{\mathbb E}\newcommand{\Prob}{\mathbb P}
\newcommand{\dd}{\,\mathrm d}\newcommand{\one}{\mathbf 1}
\DeclareMathOperator{\dist}{dist}\DeclareMathOperator{\covol}{covol}
\DeclareMathOperator{\area}{area}\DeclareMathOperator{\supp}{supp}
\numberwithin{equation}{section}
\title[Scalar collision laws]{Scalar collision laws and recognition of periodic dispersing billiards}
\author{Qian Qi}\date{October 3, 2026}
\subjclass[2020]{37D50, 52A27, 52C23, 62G05}
\keywords{Collision probabilities, reversal balance, controlled launches,
finite-range inversion, periodic billiards, repeated motifs}
\begin{document}\raggedbottom
\begin{abstract}
We recover periodic dispersing billiards from controlled scalar collision
probabilities without recording impact positions, collision angles or
collision times. At one fixed horizon, opposed launches with translated
starting distributions measure the finite difference of the obstacle
indicator. A separation and diameter bound supplies a zero on each
bounded translation chain; a minimum of partial sums then inverts that
difference exactly. The same formula recovers a smoothed occupation
function from finitely many scalar means with an error independent of
the smoothing scale. On a class with a bounded, possibly nonprimitive
period presentation, this determines the entire obstacle union and its
full translation lattice, allowing repeated shapes and arbitrary body
symmetries. On uniformly smooth classes with a specified positive
nonperiod patch-defect margin, a spatial grid of localized launch laws
gives a conditional $C^2$ reconstruction and exact discrete period
relations with confidence $1-\delta$ using
$O(\nu^{-3}\log(C/(\nu\delta)))$ attempts. The launch localization is
$O(\nu^{3/2})$; the constants and admissible accuracy depend on the
patch margin and geometric priors. This is active scalar acquisition,
not a passive or uniformly prepared count invariant. Earlier spatial
first-impact and intrinsic return-law results are preserved in full
as supplementary material.
\end{abstract}\maketitle
\section{Reversal balance and a finite-range inverse}
\label{sec:reversal}
A spatial first-impact law may reveal the boundary directly through its
support. We ask instead what remains when the output of each attempt is
only a collision bit. The starting distribution is now a controlled
input. The distinction matters: reducing the recorded output while
allowing localized preparations does not establish dominance over a
uniform-launch experiment. Our principal identity separates the two
issues. It expresses an occupation difference by two scalar collision
probabilities at the same positive horizon, with no small-time limit.

Let $\mathcal O\subset\R^d$ be closed, and put $\chi=\one_{\mathcal O}$.
An attempt starts at $q$ with velocity $v$, $|v|=1$, and runs until time
$t>0$. A start in $\mathcal O$ is an unsuccessful preparation and is
assigned output zero, as is a free start with no collision by time $t$.
A free start with a collision by that time has output one. The two zero
outcomes are not distinguished by the observation. The event includes
the endpoint time; boundary conventions are immaterial for the smooth
classes and absolutely continuous preparations below. Write $a=tv$ and
\begin{equation}\label{eq:bit}
 B_a(q)=\one_{\{q\notin\mathcal O,\ [q,q+a]\cap\mathcal O\ne\varnothing\}}.
\end{equation}
The straight segment in this definition is the unreflected path. It
detects the physical first collision, because the trajectory is straight
until that collision; no part of the path after the first collision is
observed. For a nonnegative probability density $f$ set
\[
 P_a(f)=\int_{\R^d}f(q)B_a(q)\dd q,\qquad f_a(q)=f(q-a).
\]
The denominator is the number of attempted preparations, not the number
of free or successful starts. The experiment commanding $f_a$ shifts
the starting distribution by $+a$ and reverses the velocity. It does
not run a time-reversed reflected trajectory.

\begin{theorem}[Reversal balance]\label{thm:balance}
For every closed obstacle set and every displacement $a$,
\begin{align}
 B_a(q)-B_{-a}(q+a)&=\chi(q+a)-\chi(q),\label{eq:balance-point}\\
 P_a(f)-P_{-a}(f_a)&=
       \int f(q)\{\chi(q+a)-\chi(q)\}\dd q.\label{eq:balance}
\end{align}
The pointwise identity requires no convexity, separation or periodicity.
It remains valid when the segment intersects several components.
\end{theorem}
\begin{proof}
Let $I(q)$ indicate whether $[q,q+a]$ meets $\mathcal O$. The two
oriented segments have the same image, so the left side of
\eqref{eq:balance-point} is
$\{\chi(q+a)-\chi(q)\}I(q)$. If either endpoint is in $\mathcal O$,
then $I(q)=1$; if neither is, the prefactor is zero. This proves the
pointwise assertion. In the reverse probability change variables from
$q+a$ to $q$. The result is \eqref{eq:balance}.
\end{proof}

In particular, the pair of functionals $P_a,P_{-a}$ on smooth compactly
supported launch densities determines the signed measure with density
$\delta_a\chi(q):=\chi(q+a)-\chi(q)$. Testing all nonnegative smooth
densities determines this measure uniquely, by scaling and linearity.
These are functionals of a controllable spatial input, not two real
observations. For finite acquisition we use an explicit finite family
of inputs rather than appeal to exact differentiation of such a measure.

A difference operator ordinarily leaves an additive constant along each
translation orbit undetermined. The next theorem fixes it by a geometric
zero condition, without providing a free reference point.

\begin{theorem}[Finite-chain inversion]\label{thm:chain}
Let $u\ge0$ be a function, $a\ne0$, and $m\ge1$. Suppose that for every
$x$ in a region $U$,
\begin{equation}\label{eq:zero-chain}
                 \min_{0\le k\le m}u(x+ka)=0.
\end{equation}
Then the values of $\delta_a u$ on
$U^+=\bigcup_{j=0}^{m-1}(U+ja)$ determine $u$ on $U$ by
\begin{equation}\label{eq:min-inverse}
 u(x)=-\min_{0\le k\le m}\sum_{j=0}^{k-1}\delta_a u(x+ja),
\end{equation}
where the empty sum is zero. If $u'$ obeys the same zero condition,
then, for $1\le p\le\infty$,
\begin{equation}\label{eq:chain-stability}
 \|u-u'\|_{L^p(U)}\le
       m\|\delta_a u-\delta_a u'\|_{L^p(U^+)}.
\end{equation}
More generally, perturbing every difference in a chain by at most $b$
perturbs the reconstruction by at most $mb$.
\end{theorem}
\begin{proof}
The partial sum is $u(x+ka)-u(x)$. Its minimum equals $-u(x)$ by
\eqref{eq:zero-chain}, proving the formula. The minimum of two finite
lists differs by at most their maximum entrywise difference. Hence
\[
 |u(x)-u'(x)|\le
   \sum_{j=0}^{m-1}|\delta_a u(x+ja)-\delta_a u'(x+ja)|.
\]
Minkowski's inequality and translation of each integration domain prove
\eqref{eq:chain-stability}; the uniform perturbation statement follows
from the same inequality. Almost-everywhere hypotheses suffice: only
finitely many translates of exceptional null sets occur.
\end{proof}

\begin{lemma}[A geometric zero witness]\label{lem:zero}
Suppose $\mathcal O$ is a union of compact components of diameter at
most $D_0$, with mutual distance at least $d_0>0$. If
$0<t<d_0$, $a=tv$ and $mt>D_0$, then $u=\chi$ satisfies
\eqref{eq:zero-chain} everywhere.
Fix $\sigma_0>0$ with $t<d_0-2\sigma_0$ and
$mt>D_0+2\sigma_0$. If $\kappa\ge0$ has integral one and support
in the closed unit ball, then for every $0<\sigma\le\sigma_0$ the
local occupation average
\begin{equation}\label{eq:occupation}
 u_\sigma(x)=\int \sigma^{-d}\kappa((q-x)/\sigma)\chi(q)\dd q
\end{equation}
also satisfies \eqref{eq:zero-chain}.
\end{lemma}
\begin{proof}
If all chain points belong to $\mathcal O$, successive points cannot
belong to different components, since their distance is $t<d_0$.
They therefore lie in one component, contradicting the endpoint distance
$mt>D_0$. For the second assertion, $u_\sigma(x)>0$ implies
$x\in\mathcal O+\sigma\overline B$. These enlarged components have
diameter at most $D_0+2\sigma$ and separation at least $d_0-2\sigma$.
The same argument excludes positivity at every point of the chain.
\end{proof}

\begin{corollary}[Local rigidity from scalar launch laws]\label{cor:local}
Under the diameter and separation hypotheses, the two probability
functionals at one fixed horizon and opposite directions determine the
obstacle indicator almost everywhere on $U$, using only launch densities
supported on its bounded chain enlargement. For regular closed bodies
they determine the bodies themselves in the interior of $U$.
Alternatively, the same determination follows from localized kernel
preparations at all sufficiently small scales, with the explicit
finite-chain formula at each scale.
\end{corollary}
\begin{proof}
Apply Theorem~\ref{thm:balance} to recover $\delta_a\chi$ and
Theorem~\ref{thm:chain} with Lemma~\ref{lem:zero}. Two different regular
closed convex unions would differ on an open set, so almost-everywhere
equality determines the set. For the alternative, the balance identity
recovers $\delta_a u_\sigma$ exactly, the minimum formula recovers
$u_\sigma$, and approximate-identity convergence gives $\chi$ almost
everywhere as $\sigma\downarrow0$.
\end{proof}

The zero witness is substantive. For constant sequences both zero and
one have zero difference, and for unbounded slabs parallel to $a$ the
same ambiguity can persist. Here the witness is implied by bounds on
physical component size and separation, not by an unobserved boundary
value. The local theorem uses no periodicity or boundary smoothness.
In the convex billiard class and at $t<d_0$, there cannot be a second
collision: a reflected exterior ray cannot re-enter the same convex
body, and another body is at least $d_0$ away. Thus the successful
output is also the one-collision count at this horizon, with solid
preparations still counted as zero attempts.

\section{Finite scalar acquisition of a complete patch}
\label{sec:acquisition}
We now work in the plane. Fix once and for all a nonnegative, even
$C^\infty$ probability kernel $\kappa$ supported in the unit disk, and
put $K=\|\nabla\kappa\|_1$. We may increase $K$ to be at least one.
Take, for example,
\begin{equation}\label{eq:design}
 t=\min(1,d_0/4),\qquad \sigma_0=\min(1,d_0/8),\qquad
 m=\left\lfloor\frac{D_0+2\sigma_0}{t}\right\rfloor+1,
 \qquad a=t(1,0).
\end{equation}
The chain is fixed by the numerical priors, not by an unknown boundary.
For $f_{x,\sigma}(q)=\sigma^{-2}\kappa((q-x)/\sigma)$ define
\[
 d_\sigma(x)=P_a(f_{x,\sigma})-P_{-a}(f_{x+a,\sigma}).
\]
Equation~\eqref{eq:balance} and Lemma~\ref{lem:zero} give
\begin{equation}\label{eq:smoothed-min}
 d_\sigma(x)=u_\sigma(x+a)-u_\sigma(x),\qquad
 u_\sigma(x)=-\min_{0\le k\le m}\sum_{j<k}d_\sigma(x+ja).
\end{equation}
There is no deconvolution, time derivative or noise amplification by
$\sigma^{-1}$ in this step.

\begin{proposition}[Stable recovery of occupation averages]
\label{prop:mean-stability}
At each point $x$ the $2m$ scalar means in \eqref{eq:smoothed-min}
determine $u_\sigma(x)$. If each mean has absolute error at most
$\epsilon$, the error is at most $2m\epsilon$, uniformly in $\sigma$.
Clipping the answer to $[0,1]$ cannot increase this error. For two
physical candidates whose corresponding means differ by at most $b$,
their occupation averages at $x$ differ by at most $2mb$.
\end{proposition}
\begin{proof}
Each balance difference has error at most $2\epsilon$, or at most
$2b$ in the comparison version. Apply the perturbation conclusion of
Theorem~\ref{thm:chain}. The true average lies in $[0,1]$.
\end{proof}

Each mean is estimated by repeated attempts with the prescribed input
law. The actual start position need not be recorded with the output;
the commanded distribution, its center and the direction are known.
The only per-attempt recorded variable is the bit. If the actual launch
law differs from its command by total variation at most $\zeta$, the
bias of every mean is at most $\zeta$. In particular a center error
$\ell$ changes this smooth input density in $L^1$ by at most
$K\ell/\sigma$. These observations describe input calibration costs;
they are not a claim that arbitrary directional or dynamical apparatus
errors are harmless.

\begin{lemma}[A finite grid controls the patch geometry]
\label{lem:grid}
Assume all physical bodies are strictly convex, have curvature at most
$\kappa_+$, and satisfy the common separation and diameter bounds.
Let $U$ be a square containing, with a collar of width $2D_0+2$, every
body to be compared. Take
$\sigma<\min(\sigma_0,(2\kappa_+)^{-1},d_0/20,1/4)$.
A grid of covering radius $\sigma/(32K)$ in $U$ contains
$O_U(\sigma^{-2})$ points. If the scalar means for two admissible
physical candidates differ at every grid command by at most
$b=1/(32m)$, then their complete bodies in the protected inner patch
can be matched with Hausdorff error at most $2\sigma$.
\end{lemma}
\begin{proof}
By Proposition~\ref{prop:mean-stability} the occupation averages differ
by at most $1/16$ at each grid point. Each occupation function has
Lipschitz constant at most $K/\sigma$. The chosen covering radius
therefore gives a uniform difference at most $1/8$ on $U$.

For a convex body $C$, let $C_{-\sigma}$ denote its erosion by the
closed radius-$\sigma$ disk. If $x\in C_{-\sigma}$, then
$u_\sigma(x)=1$, whereas if $\dist(x,\mathcal O')>\sigma$, the other
occupation average is zero. It follows that
\begin{equation}\label{eq:erosion-inclusion}
 \mathcal O_{-\sigma}\cap U\subset\mathcal O'+\sigma\overline B,
 \qquad
 \mathcal O'_{-\sigma}\cap U\subset\mathcal O+\sigma\overline B.
\end{equation}
The notation on the left means the union of the component erosions.
The curvature upper bound ensures an interior rolling radius
$r=1/\kappa_+$. One elementary justification uses the support function:
$p+p''\ge r$, so $p-r$ is a convex support function and
$C=C_0+r\overline B$. Hence
$C=(C_{-\sigma})+\sigma\overline B$ when $\sigma<r$.
Each erosion is nonempty and connected.

Distinct expanded target components have positive separation. A
connected erosion in \eqref{eq:erosion-inclusion} must therefore lie
in the radius-$\sigma$ enlargement of a single target body. Restoring
the disk shows $C\subset D+2\sigma\overline B$. Apply the other
inclusion to $D$ to obtain $D\subset E+2\sigma\overline B$ for a
single source body $E$. All these bodies and their erosions are in $U$
by the protective collar. Since $C\subset E+4\sigma\overline B$ and
$4\sigma<d_0$, necessarily $E=C$. Thus $d_H(C,D)\le2\sigma$.
The same argument in the reverse direction makes the matching bijective
for all complete bodies used in the protected patch. No incomplete
outer fragment is declared to be a physical component.
\end{proof}

For later use we record the smooth interpolation underlying the
$C^2$ conclusion. Here $C^j$ norms of support functions are on the
normal-angle circle, in a fixed laboratory frame.
\begin{lemma}[From Hausdorff error to smooth support error]
\label{lem:smoothing}
For convex bodies with bounded $C^6$ support norms, Hausdorff distance
at most $e\le1$ implies $C^2$ support difference at most $C e^{2/3}$.
The same estimate holds for centered support functions, with a changed
constant and their Steiner-center error bounded by $2e$.
\end{lemma}
\begin{proof}
Let $\varphi$ be a smooth even probability kernel of compact support
on the line, and set $\varphi_2(s)=\varphi(s/2)/2$ and
$J=(4\varphi-\varphi_2)/3$. Then $J$ has integral one and its moments
of degrees one, two and three vanish. Periodized convolution at scale
$h$ satisfies
\[
 \|J_h*(p-p')\|_{C^2}\le Ceh^{-2},\qquad
 \|J_h*p-p\|_{C^2}+\|J_h*p'-p'\|_{C^2}\le C h^4.
\]
The first inequality uses $\|p-p'\|_\infty=d_H(C,D)$ and the $L^1$
norms of the first two derivatives of the kernel. The second is Taylor's
formula for $p,p',p''$ through the cancelled orders, using the $C^6$
bound. Choose $h=e^{1/6}$. The Steiner point is
$\pi^{-1}\int_0^{2\pi}p(\theta)(\cos\theta,\sin\theta)\dd\theta$;
its change has norm at most $2e$. Subtracting this translation gives
the centered estimate.
\end{proof}

An admissible-candidate formulation will suffice below. It avoids
claiming an efficient reconstruction algorithm for all smooth shapes:
we select a physical candidate consistent with the finitely many scalar
means, then apply the explicit finite period tests to its recovered
patch. The statements concern both this selection and its quantified
error, not the execution time of searching an infinite-dimensional class.

\section{Recognition of the intrinsic period lattice}
\label{sec:periods}
The scalar inverse recovers a physical patch before periodicity is used.
We next give the local period argument needed to extend that patch.
This argument is retained from the preceding repeated-motif revision
\cite{Supplement}; it is included here so that the primary theorem does
not depend on a different sensor's acquisition proof.

\begin{definition}[Bounded periodic class]\label{def:class}
Fix $r_0,L_0,V_0,D_0,d_0,\kappa_-,\kappa_+>0$. The class
$\mathcal B$ consists of nonempty planar unions
\begin{equation}\label{eq:period-presentation}
 \mathcal O=\bigcup_{i=1}^r(C_i+\Lambda_0),\qquad
 \Lambda_0=L\Z^2,\qquad 1\le r\le r_0,
\end{equation}
with exactly $r$ component orbits modulo $\Lambda_0$,
$\|L\|,\|L^{-1}\|\le L_0$, and $\covol\Lambda_0\ge V_0$.
Every physical component is a compact strictly convex body of diameter
at most $D_0$, with curvature in $[\kappa_-,\kappa_+]$; distinct
components are separated by at least $d_0$. Centered support functions
have $C^{6,\beta}$ norm at most $M_6$, for a fixed $0<\beta\le1$.
The unknown presentation need not be primitive. Repeated shapes and
arbitrary symmetries are allowed. Its full period group is
$\Pi(\mathcal O)=\{w:\mathcal O+w=\mathcal O\}$.
\end{definition}

The high-order bound will only enter finite $C^2$ reconstruction.
Set $B=1+L_0$ and $Q_s=[-s,s]^2$. A convenient fixed observation and
launch aperture is
\begin{equation}\label{eq:aperture}
 R_0=12B+3D_0+4,\qquad R=R_0+mt+2,\qquad Q_R,
\end{equation}
where $t,m$ are from \eqref{eq:design}. Scalar preparations centered
on $Q_{R_0}$ and on the forward chains needed for it are supported in
$Q_R$ for every $\sigma\le\sigma_0$. Each physical flight has the same
fixed duration $t$. There is no recentering at an estimated obstacle.

For a component $C$ write $z_C$ for its Steiner point and $p_C$ for its
centered support function. In the laboratory directions define
\begin{align}
 d_*(C,D)&=\|z_C-z_D\|_\infty+\|p_C-p_D\|_\infty,\label{eq:bodymetric}\\
 \mathcal D_{\mathcal O}(w)&=
 \max_{C:z_C\in Q_B}\ \max_{s\in\{-1,1\}}
                   \inf_{D\in\mathcal C}d_*(C+sw,D),\label{eq:defect}
\end{align}
where $\mathcal C$ is the set of physical components. The metric does
not optimize over rotations. Hausdorff error $e$ implies center error
at most $2e$ and centered-support error at most $3e$. All relevant
infima are attained, since distant centers cannot minimize them.

\begin{lemma}[A finite test for a global period]\label{lem:period}
One has $\mathcal D_{\mathcal O}(w)=0$ if and only if
$w\in\Pi(\mathcal O)$. For $\|w\|_\infty\le8B$ the zero test uses
only components with centers in $Q_B$ and $Q_{9B}$.
\end{lemma}
\begin{proof}
Every $\Lambda_0$-component orbit has a representative with center in
$L[-1/2,1/2]^2\subset Q_B$. If the defect vanishes, each such
representative has both translates by $+w$ and $-w$ among the physical
components. Translating by $\Lambda_0$ gives
$\mathcal O+w\subset\mathcal O$ and $\mathcal O-w\subset\mathcal O$.
The latter implies the reverse of the former inclusion. The converse
is immediate. For the bounded zero test, every possible equal target
has center $z_C\pm w\in Q_{9B}$.
\end{proof}

\begin{lemma}[A protected generating list]\label{lem:generators}
The group $\Pi=\Pi(\mathcal O)$ is a lattice containing $\Lambda_0$,
\begin{equation}\label{eq:index}
 [\Pi:\Lambda_0]\le r\le r_0,\qquad \covol\Pi\ge V_*:=V_0/r_0.
\end{equation}
For any root component centered in $Q_{2B}$, take its center differences
to components centered in $Q_{4B}$ and retain those passing
Lemma~\ref{lem:period}. Their integer span is exactly $\Pi$.
\end{lemma}
\begin{proof}
Map a period coset modulo $\Lambda_0$ to the $\Lambda_0$-orbit of the
translated root. This is injective, since a compact body has no nonzero
translation symmetry. There are at most $r$ such cosets. A finite
extension of a full-rank lattice is a full-rank lattice; the covolume
formula proves \eqref{eq:index}. This quotient group acts freely on
the $r$ component orbits, so its order divides $r$.

All accepted differences are true periods. Conversely, the two columns
of $L$ and representatives of all cosets of $\Pi/\Lambda_0$ can be
chosen to have norm less than $B$, using the same centered fundamental
parallelogram. Translating the root by any of these vectors gives a
target in $Q_{3B}$, strictly inside the cutoff $Q_{4B}$. The accepted
list contains generators of $\Lambda_0$ and representatives of every
coset, and hence generates $\Pi$.
\end{proof}

\begin{theorem}[Whole-table rigidity from scalar collision laws]
\label{thm:global-exact}
On $\mathcal B$, the collision-probability functionals for the two
opposed directions at the fixed duration \eqref{eq:design}, with
controlled smooth starting densities in $Q_R$, determine the entire
obstacle union, its full period lattice $\Pi$, its primitive number
$r_*$ of component orbits, and primitive free area $A_*$. The same is
true for the family of localized kernel commands at all sufficiently
small scales. No impact position, impact angle, collision time, shape
label, integer copy label, primitive cell or free-area normalization
is an observation. The laboratory frame is supplied by the launch
controls; forgetting it leaves a common Euclidean-isometry ambiguity.
\end{theorem}
\begin{proof}
Corollary~\ref{cor:local} recovers the occupation indicator on
$Q_{R_0}$, and regular closedness identifies the physical set there.
Every component with center in $Q_{10B}$, and every component needed
to recognize one of these, lies wholly within the protective collar
of this square. Physical separation identifies the complete components.
Lemmas~\ref{lem:period} and \ref{lem:generators} then recover $\Pi$.

All $\Pi$-component orbits have representatives in $Q_B$. Among
components centered in $Q_{2B}$, two are in the same orbit exactly when
their center difference passes the period test. Indeed a period maps
the first to a physical component with the second's Steiner point;
two disjoint full-dimensional convex bodies cannot have the same
Steiner point, which lies in their interiors. Thus the relation gives
$r_*$ and one representative per orbit. These and $\Pi$ determine the
union everywhere. Finally
\begin{equation}\label{eq:area}
 A_*=\covol\Pi-\sum_{j=1}^{r_*}\area(C_j).
\end{equation}
The conclusion concerns the intrinsic union, not a chosen multiple-cell
presentation. Basis computation from finite exact period vectors is
justified quantitatively in Lemma~\ref{lem:locking} below.
\end{proof}

\begin{definition}[Nonperiod patch margin]\label{def:margin}
For $\eta>0$, let $\mathcal B_\eta$ consist of those members of
$\mathcal B$ such that
\begin{equation}\label{eq:margin}
 w=z_D-z_C,\quad z_C\in Q_{3B},\ z_D\in Q_{5B},\quad
 w\notin\Pi(\mathcal O)\quad\Longrightarrow\quad
                  \mathcal D_{\mathcal O}(w)\ge\eta.
\end{equation}
This concerns translations of a whole central patch, not separation
between individual body shapes. For each fixed table it holds with some
positive $\eta$, because there are finitely many candidate pairs and
Lemma~\ref{lem:period} makes every nonperiod defect positive.
\end{definition}

The next deterministic statement is the link between geometric estimation
and exact discrete information. It does not suppose that perturbed real
vectors themselves generate a discrete group.
\begin{lemma}[Robust period classification and rational locking]
\label{lem:locking}
Fix the numerical bounds of $\mathcal B_\eta$. Suppose complete
components on the protected patch are given by matched approximations
of Hausdorff error at most $e$. For sufficiently small $e$, depending
on $\eta$ and the geometric bounds, finite comparisons recover the true
primitive orbit count and the exact rational coordinates of an accepted
period-generating list in an independent-pair basis. They give a matched
basis for $\Pi$ with error $O(e)$ and an area estimate with error $O(e)$.
\end{lemma}
\begin{proof}
Use approximate centers in $Q_{2B}$ as sources, choose a root there,
and use its targets in $Q_{4B}$. A measured difference $\widehat w$
has a corresponding true center difference $w$ with error at most
$4e$. Test both signs of this translation on all selected sources,
minimizing $d_*$ over approximate components centered in $Q_{10B}$.
For each source-target comparison, the center term changes by at most
$8e$ and the centered-support term by at most $6e$.

If $w$ is a period, all needed partners lie within the target cutoff
by the strict aperture margins, and the measured defect is at most
$14e$. If it is not, its true endpoints lie in $Q_{3B},Q_{5B}$, and
\eqref{eq:margin} applies. Every witness with true center in $Q_B$ is
selected. Restricting the targets cannot reduce a true positive infimum,
so the measured defect is at least $\eta-14e$. Reserve another $6e$
for numerical enclosures of centers and support suprema. Threshold
$\eta/2$, with $20e<\eta/4$, accepts exactly the true periods among
these candidates. The same comparison on all pairs of selected sources
gives their true orbit equivalence relation and hence $r_*$. Targets
for these tests also lie within $Q_{10B}$. Lemma~\ref{lem:generators}
and its strict cutoff margin ensure all protected generators are present.

All accepted vectors have norm at most $U=12B$ for small $e$. A nonzero
determinant of two true periods is an integer multiple of
$\covol\Pi\ge V_*$. Determinant perturbations are $O(Ue)$; after
shrinking $e$, independence is decided at threshold $V_*/2$. Choose
an independent accepted pair with true matrix $D$. The sublattice
$D\Z^2$ has index
\[
       n=|\det D|/\covol\Pi\le Q:=\max(1,\lceil U^2/V_*\rceil).
\]
Every coordinate of $D^{-1}w_j$ has reduced denominator dividing $n$,
and a prior-bounded numerator. The inverse-matrix identity bounds its
estimated error by $C_*e$. Distinct rationals of denominator at most
$Q$ differ by at least $Q^{-2}$. If $C_*e<1/(3Q^2)$, finite search
therefore identifies each exact coordinate uniquely.

Let $q_*$ be their common denominator; it divides $n$ and is at most
$Q$. A column Hermite basis $H$ of the integer span of $q_*I_2$ and
$q_*D^{-1}w_j$ gives
\begin{equation}\label{eq:hermite}
                  \Pi=(DH/q_*)\Z^2.
\end{equation}
The integer group contains $q_*\Z^2$, so the Hermite entries are
bounded in terms of $Q$. Replacing $D$ by its measured version gives
a matched basis with error $O(e)$. This basis need not be a continuously
chosen reduced basis. The parallel-body bound gives an $O(e)$ area
error for each matched convex component; use \eqref{eq:area} for
$A_*$. Certified finite angular meshes and the uniform support Lipschitz
bound enclose all suprema and center integrals to the reserved precision.
\end{proof}

\section{Finite confidence from localized scalar preparations}
\label{sec:finite}
The constants below may depend on the fixed kernel and every numerical
bound of $\mathcal B_\eta$, including $\eta$. We count independent
attempted launches. Spatial preparation precision, deterministic search,
arithmetic and apparatus travel are different resources.

\begin{theorem}[A finite scalar experiment on a patch-margin class]
\label{thm:finite}
On $\mathcal B_\eta$ there are $C,c,\nu_0>0$ such that, for
$0<\nu<\nu_0$ and $0<\delta<1/4$, a deterministic finite list of
localized launch distributions in the fixed square $Q_R$, with only
the two directions $\pm(1,0)$ and the fixed duration $t$, yields the
following guarantee. With probability at least $1-\delta$, it recovers
the primitive orbit count and exact rational relations of a generating
list for $\Pi$, and gives a finite presentation with matched $C^2$
support error, matched period-basis error and primitive-area error at
most $C\nu$. The list uses localization scale
\begin{equation}\label{eq:scalar-scale}
                         \sigma=c\nu^{3/2}
\end{equation}
and at most
\begin{equation}\label{eq:scalar-cost}
                    C\nu^{-3}\log\frac{C}{\nu\delta}
\end{equation}
independent attempts, including solid starts and misses. Only one bit
is recorded per attempt. No output collision localization is required.
A sufficiently small fixed bound on each command's probability bias
is allowed, as specified in the proof.

The uniform exact discrete decisions require the known patch margin
$\eta$. The rate is a conditional sufficient upper bound, not a minimax
rate or an end-to-end computation bound. Localized spatial preparation
and collimated launch directions are substantive inputs.
\end{theorem}
\begin{proof}
\emph{A finite deterministic experiment.}
Take the grid of Lemma~\ref{lem:grid} on $Q_{R_0}$, together with
all its commanded chain centers. Its number of scalar means is
$J\le C m\sigma^{-2}$. At each command estimate its Bernoulli mean
to statistical error at most $\epsilon_0/2$, where
$\epsilon_0=1/(256m)$ is fixed. Allow a systematic mean bias at most
$\epsilon_0/2$. Then each reported mean is within $\epsilon_0$ of
the ideal mean. Hoeffding's inequality and a union bound make this
simultaneous event have probability at least $1-\delta$ with
$C\epsilon_0^{-2}\log(2J/\delta)$ independent attempts per command.
The experiment is deterministic before sampling. It uses no separate
pilot or estimated onset. The total is
\begin{equation}\label{eq:raw-cost}
                      C\sigma^{-2}\log(C/(\sigma\delta)).
\end{equation}

\emph{Selection of an admissible patch.}
The bounded periodic class is precompact in local $C^2$ geometry, and
may be replaced by its closure in a weaker smooth topology retaining
the displayed bounds. To see this, choose the bounded presentations,
pass to a subsequence with fixed $r$, choose component centers in a
closed bounded fundamental parallelogram, and use compactness of the
matrices and support functions. Positive separation and curvature
persist. Higher derivative bounds persist in their usual weak or
H\"older-compact sense. This description also gives a separable
parameter space; a closed compact enlargement satisfying the same
bounds suffices throughout. We do not require the patch-margin subclass
to be closed, since the candidates may be selected in $\mathcal B$.

For each of the finitely many commands the ideal probability is
continuous on this physical class. Except on a null set of initial
points, whether a free segment first meets a strictly convex component
before its endpoint is unchanged under sufficiently small geometric
perturbations. The excluded points lie on body boundaries, translated
endpoint boundaries or tangent lines. Only finitely many bodies can
meet the bounded launch-and-flight region. Dominated convergence proves
continuity of the integral against the fixed smooth command density.

Minimize the maximum mean discrepancy over this class. One can choose
a Borel approximate minimizer within $\epsilon_0$ of the infimum from
a fixed countable dense family: order the family and take its first
eligible member. The infimum of countably many continuous functions
of the data is measurable, so the rule is measurable. On the simultaneous
mean event the true table has discrepancy at most $\epsilon_0$;
the chosen candidate has discrepancy at most $2\epsilon_0$. Its ideal
means differ from those of the truth by at most $3\epsilon_0<1/(32m)$.
Lemma~\ref{lem:grid} gives matched Hausdorff error $e=2\sigma$ for
all complete components used in the central period tests. This step
is an existence-based regularization, not a polynomial-time search
algorithm for arbitrary smooth obstacles.

\emph{Discrete and smooth recovery.}
Use those candidate components as the geometric approximations in
Lemma~\ref{lem:locking}, applying the known margin of the true table.
The candidate itself need not have the same exact symmetry group:
period tests are approximate tests against the true margin, not a
claim that the candidate's raw periods are correct. When $e$ is below
the thresholds in that lemma, the recovered orbit count and rational
relations are the true ones and the matched lattice basis error is
$O(e)$. Choose one candidate body per recovered true orbit and repeat
these bodies under the reconstructed basis. Lemma~\ref{lem:smoothing}
gives $C^2$ support error $O(\sigma^{2/3})$, including the center error.
The basis and its inverse remain bounded. Consequently only a bounded
set of integer translates can enter any bounded fundamental region;
positive physical separation persists for small $\sigma$ and the
new presentation defines a valid periodic union. Its primitive free
area error is $O(\sigma)$ by \eqref{eq:area} and the area estimate.

Choose \eqref{eq:scalar-scale} with a small constant and decrease
$\nu_0$ to impose the rolling-radius, separation, patch-gap and rational
thresholds. This proves the geometric accuracy and transforms
\eqref{eq:raw-cost} into \eqref{eq:scalar-cost}. No additional
independence is needed for root or candidate selection, since every
conclusion holds on the one simultaneous deterministic-command event.
\end{proof}

The finite-presentation metric just used compares one matched body for
each true primitive component orbit, their support functions in the
common laboratory directions, and matched period bases. It does not
assert bounded Hausdorff distance between two entire infinite unions
whose periods differ slightly. Such a metric would encode a different
question. The exact rational relations are relations within the chosen
period-generating list; they are not exact Euclidean coordinate values.

\begin{remark}[Calibration and the meaning of the rate]\label{rem:resources}
A command implemented within total variation $\epsilon_0/2$ obeys the
bias allowance. With this kernel, spatial center error at most
$c_\kappa\epsilon_0\sigma$ is one sufficient contribution to that
allowance. The theorem supposes the two directions and the horizon
are calibrated; it does not prove a uniform angular-jitter bound near
grazing rays. Only their combined certified probability bias, when
available, may be charged to the same allowance. The output bit does
not certify input accuracy. Increasing spatial resolution therefore
still costs preparation precision, even though a fixed absolute accuracy
for each scalar mean suffices. This cost, the number of commands,
attempt count, bit arithmetic, and the infinite-class selection problem
are kept separate.
\end{remark}

\begin{corollary}[Unknown patch margin]\label{cor:pointwise}
For each fixed member of $\mathcal B$, a sequence of fresh finite
scalar experiments with $\sigma_j\downarrow0$ and summable simultaneous
failure probabilities gives convergent local geometry and eventually
correct primitive orbit count and rational period relations almost
surely. A threshold tending to zero more slowly than $\sigma_j$, such
as $\sqrt{\sigma_j}$, suffices for the patch comparisons. No uniform
finite stopping certificate over all of $\mathcal B$ is asserted.
\end{corollary}
\begin{proof}
Every relevant false translation in a fixed bounded patch has positive
defect, and there are finitely many such candidates. Let their minimum
be $\eta_{\mathcal O}>0$. The geometric errors are $O(\sigma_j)$,
so the proposed threshold eventually lies strictly between true and
false defects. All protected generators remain present, and the fixed
rational thresholds eventually hold. By the first Borel--Cantelli
lemma, only finitely many simultaneous mean events fail. Recompute
rather than permanently retain earlier classifications. The geometric
error estimate converges and the discrete decisions are eventually
correct. The unknown $\eta_{\mathcal O}$ does not supply an observable
finite index at which this has occurred.
\end{proof}

\begin{proposition}[Repeated motifs near a symmetry increase]
\label{prop:jump}
The patch-margin dependence cannot be removed from a uniform continuous
recovery assertion for the full period lattice and primitive orbit count
on the whole class $\mathcal B$.
\end{proposition}
\begin{proof}
Use the retained repeated-disk family \cite{Supplement}: disks of
radius $1/4$ centered at $(4k,4l)$ and $(4k+2+s,4l)$, where
$0\le s<1/8$ and $k,l\in\Z$. All displayed geometric and smoothness
bounds are uniform. At $s=0$ the full lattice is $2\Z\times4\Z$
and there is one component orbit; at $s>0$ it is $4\Z\times4\Z$
and there are two. Indeed a horizontal exchange translation preserves
$\{0,2+s\}$ modulo $4$ exactly when $2(2+s)$ is divisible by $4$.
The exchange candidate has central defect $2s$ for $s>0$. It can be
arbitrarily close to a period while not being one. The bodies converge
smoothly on every bounded patch but the exact discrete invariants jump.
This geometric example applies to any reconstruction derived from close
patch geometry. It does not contradict exact determination or the
pointwise conclusion above.
\end{proof}

\section{Information categories, local periodicity and retained results}
\label{sec:comparison}
\subsection{What is observed and what is reconstructed}
There are three different collision experiments in the present submission
package. The spatial first-impact experiment of Supplement P launches
uniformly in a fixed laboratory square with independent uniform directions
and records a localized first-impact position. In exact data, the support
of its position law is the relevant physical boundary union. Its protected
launch collars give positive mass to every boundary arc. The exact
finite-patch inverse in that experiment is therefore informationally
direct; it is not boundary recovery from hidden travel-time or spectral
data.

The experiment of this primary article records no collision coordinates.
It instead allows localized commanded launch densities and two opposed
collimated directions. The measured imbalance is a finite difference
of occupation, not the occupation function itself. The cancellation in
Theorem~\ref{thm:balance} removes paths which enter and leave a body
between two free endpoints. The finite-chain zero condition restores
the additive information lost in taking that difference. The same
nonlinear inverse works for mollified occupation without differentiating
noisy probability functions. The theorem is thus a scalar-output
alternative within the collision-law programme, with an explicitly
richer preparation contract than a single uniform launch law. Neither
experiment is identified with the other by an unsupported information
ordering.

Earlier intrinsic return-law results concern different selected
itineraries, endpoint arclength laws and moment compressions. Their
count-only finite fibers, formal fibers and smooth Taylor agreements
remain distinct from equality of exact analytic count germs. A family
of spatially commanded binary collision probabilities is not the
previously studied uniformly prepared count-only datum. The present
argument makes no inference about that analytic count-germ problem,
passive trajectories, marked-length spectra or spectral rigidity.

\subsection{Finite differences and local periodicity}
The finite-difference operator in \eqref{eq:balance} is elementary.
The point is the physical identity realizing it without output positions,
and the bounded zero-witness inverse, rather than a new general theorem
about difference operators. Covariogram formulas also link translations
to occupation statistics: for a measurable finite-volume set, the
covariogram integrates the product of two translated indicators.
Galerne~\cite{Galerne} relates its directional derivatives to perimeter
and directional variation. Here spatially specified test densities
retain the location of a \emph{signed difference}; the datum is not a
translation-invariant autocorrelation. No covariogram uniqueness or
priority conclusion follows from the comparison.

There is also a distinction between recognizing a period and proving
that a set must be periodic. Lagarias and Pleasants~\cite{LP} study
patch-counting complexity and crystallinity of Delone sets. Dolbilin,
Garber, Schulte and Senechal~\cite{DGSS} study regularity radii: local
cluster equivalence can impose global regularity, with quantitative
bounds in the classes they treat. Herva and Kari~\cite{HK} develop
local-complexity and annihilator methods for Delone configurations,
including periodic decomposition and criteria forcing periodicity.
These are local-to-global structural questions with their own hypotheses.

In contrast, \eqref{eq:period-presentation} already assumes a bounded
periodic presentation. Its bound ensures a known finite central patch
contains representatives of every orbit and a protected list containing
generators of the full period group. Our period test does not infer
crystallinity for an arbitrary Delone set from observing a large patch.
The new scalar balance need not vanish and is not an annihilator
hypothesis. Its inversion first recovers the physical patch; only then
is the previously proved finite translation criterion used. This
separation places the contribution without suggesting that bounded
period recognition supersedes general local-periodicity theory.

\subsection{Two finite experiments and their different resources}
On the repeated-motif patch-margin class, the retained spatial first-hit
theorem gives a sufficient attempt bound
$C_\eta\nu^{-3/4}\log(C_\eta/(\nu\delta))$, with output localization
$O(\nu^{3/2})$, for its uniform-launch boundary-coverage design. The
scalar theorem here gives
$C_\eta\nu^{-3}\log(C_\eta/(\nu\delta))$, with input localization
$O(\nu^{3/2})$ and no position output, using a two-dimensional grid of
commands and a fixed precision per scalar mean. Both constants and
accuracy thresholds depend on the same kind of nonperiod patch margin,
in addition to their respective apparatus and geometric bounds. These
are sufficient upper bounds for different experiments, not competing
minimax rates. The former's boundary coverage and the latter's scalar
spatial grid count different operations.

The boundary smoothing estimate is a compensated-convolution argument,
and the probability estimate is an elementary bounded-variable inequality.
Neither exponent is offered as a separate assertion of global optimality.
The exact reversal identity, finite-chain inverse and the explicit
preparation/output trade are the additions; the geometric period lemma
and arithmetic classification retain their preceding provenance.

\subsection{Submission architecture and source status}
Supplement P is the entire reviewed A2 v28 paper and its nested source
history, supplied without changing any mathematical input. It includes
the repeated-motif and translation-separated first-impact theorems,
record-local collision clouds, intrinsic return-law inverses, completion,
sequential certificates, moment theory and the relative-law supplementary
volumes. Those results are neither deleted nor asserted anew by the
primary article. The present text includes the period proof it uses
and has no dependency on the older return-law chapters.

The phrase ``already published'' in the preserved v28 README refers to
a repository deposit and must not be read as journal publication of v27.
The corrected status in this submission is \emph{previously deposited
and source-qualified author revision}. A prior exact-source build
qualifies its own frozen source only; the new primary and retained
volumes need the current qualification receipt. These source statements
do not replace independent proof review or a publication record.

\subsection*{Acknowledgment and responsibility}
AI-assisted analysis contributed to the reversal identity, finite-chain
inverse and scalar-acquisition arguments and to the source checks. The
named author is responsible for verifying the proofs and references
before submission. Finite diagnostics and compilation are not independent
mathematical certification or an editorial decision.
\begin{thebibliography}{99}\raggedright
\bibitem{DGSS}
N. Dolbilin, A. Garber, E. Schulte, and M. Senechal,
\emph{Bounds for the regularity radius of Delone sets},
Discrete Comput. Geom. \textbf{74} (2025), 78--94.
\href{https://doi.org/10.1007/s00454-024-00666-6}{doi:10.1007/s00454-024-00666-6}.
\bibitem{Galerne}
B. Galerne, \emph{Computation of the perimeter of measurable sets via
their covariogram. Applications to random sets}, Image Anal. Stereol.
\textbf{30} (2011), 39--51.
\href{https://doi.org/10.5566/ias.v30.p39-51}{doi:10.5566/ias.v30.p39-51}.
\bibitem{HK}
P. Herva and J. Kari, \emph{Periodicity and local complexity of Delone
sets}, arXiv:2504.20709v1 (2025).
\bibitem{LP}
J. C. Lagarias and P. A. B. Pleasants, \emph{Local complexity of Delone
sets and crystallinity}, arXiv:math/0105088v1 (2001).
\bibitem{Supplement}
Q. Qi, \emph{Reference-free certification from intrinsic boundary laws},
A2 v28, author revision of October 3, 2026. Complete reviewed source
included as Supplement P, with its retained supplementary volumes;
not cited as a journal publication.
\end{thebibliography}

\end{document}
